\subsection{Synthetic coverage exposes sensitivity to serial dependence}
% Direct fields formatted from the audited coverage summary.
\newcommand{\CoverageReplicates}{300}
\newcommand{\CoverageTotalDatasets}{3600}
\newcommand{\CoverageTwoStageOUKnownGridChanges}{68}

All \CoverageTotalDatasets{} prespecified synthetic datasets finished without
failed draws. Each of the 84 delay-method comparisons returned an interval for
every requested dataset, so its method-defined and unconditional coverage rates
coincide. Table~\ref{tab:coverage-delay} gives exact counts for four procedures;
Fig.~\ref{fig:coverage-delay} includes all seven and their Monte Carlo uncertainty.

Under IID Gaussian generation, the IID local and profile procedures gave
coverage near the nominal reference. Applying those same procedures under OU
generation produced severe undercoverage, especially at the longer correlation
length. Estimating a common variance did not compensate for ignoring serial
covariance. HAC intervals improved coverage under correlated generation, with
substantial bandwidth sensitivity. Even the largest tested bandwidth fell below
the nominal reference in the OU60 comparisons. The synthetic thermal mean is
correct by construction here; these failures do not require between-record
mean discrepancy.
% Generated by render_coverage_table.py; see coverage_table_manifest.json.
\begin{table}[tbp]
\centering\small
\begin{tabular}{llrrrr}
\toprule
Model & Error law & IID known & IID fitted & HAC60 & Profile fitted \\
\midrule
$n=1$ & IID & 278/300 & 278/300 & 273/300 & 278/300 \\
$n=1$ & OU10 & 105/300 & 103/300 & 263/300 & 102/300 \\
$n=1$ & OU60 & 73/300 & 66/300 & 235/300 & 66/300 \\
$n=2$ & IID & 281/300 & 280/300 & 272/300 & 280/300 \\
$n=2$ & OU10 & 110/300 & 110/300 & 266/300 & 108/300 \\
$n=2$ & OU60 & 61/300 & 56/300 & 249/300 & 55/300 \\
$n=12$ & IID & 282/300 & 281/300 & 276/300 & 280/300 \\
$n=12$ & OU10 & 106/300 & 104/300 & 262/300 & 104/300 \\
$n=12$ & OU60 & 59/300 & 53/300 & 248/300 & 53/300 \\
Pure delay & IID & 290/300 & 288/300 & 283/300 & 288/300 \\
Pure delay & OU10 & 106/300 & 104/300 & 273/300 & 103/300 \\
Pure delay & OU60 & 74/300 & 68/300 & 254/300 & 68/300 \\
\bottomrule
\end{tabular}
\caption{Delay-interval coverage counts for four selected procedures. Each entry is covered/defined, from \CoverageReplicates{} requested datasets at that generating point and true covariance; -- would denote no defined interval. All intervals have nominal 95\% reference levels. IID known/fitted refer to Gaussian log-Wald intervals; HAC60 uses a Bartlett bandwidth of 60 retained samples; profile fitted uses a common profiled IID variance and retains disconnected components. OU10/OU60 are true correlation lengths in seconds, with fixed marginal standard deviation. Full method-specific failure bounds and grid diagnostics are retained in the coverage CSV; Wilson Monte Carlo intervals and all seven methods appear in Fig.~\ref{fig:coverage-delay}.}
\label{tab:coverage-delay}
\end{table}


The profile comparison retains a numerical limitation. Three-point quadratic
interpolation changed \CoverageTwoStageOUKnownGridChanges{} of
\CoverageReplicates{} coverage classifications for the two-stage known-variance
profile under OU60 generation. This is a sensitivity check on the saved grid,
without new nuisance optimization or an endpoint-accuracy certificate. Both
primary and secondary classifications remain conditional diagnostics of an IID
likelihood applied to correlated errors. The full CSV retains this diagnostic
for every profile case; interpreting the known-variance profile rates as exact
continuous-profile coverage would overstate their numerical verification.

Figure~\ref{fig:coverage-width} compares median interval widths. The one-stage
profile width sums its two accepted pole-label components rather than measuring
the gap between them. A wider accepted set can therefore concern a different
global parameter representation from a local interval around one pole label.
These simulations fix four generating points, the recorded inputs/times and the
initial-temperature, ambient and heater-calibration assumptions. They do not
measure coverage of the real-data block-bootstrap intervals or uncertainty from
the observed between-record discrepancy.

\begin{figure}[p]
 \centering
 \includegraphics[width=\linewidth]{coverage_delay_coverage.pdf}
 \caption{Serial dependence reduces coverage of IID-based delay intervals in
 controlled simulation. Each point is covered/defined over
 \CoverageReplicates{} requested datasets for one model and true error process;
 bars are pointwise 95\% Wilson intervals for Monte Carlo sampling uncertainty.
 All displayed procedures returned intervals for every dataset. The dashed line
 is nominal 0.95 coverage. Known/fitted denote fixed or profiled IID variance;
 local IID intervals are Gaussian log-Wald intervals and profiles retain all
 accepted components. HAC labels give Bartlett bandwidths in retained samples.
 All true errors have the same marginal standard deviation; rows change their
 correlation law. The thermal mean and startup are correctly specified in each
 generating scenario. Profile sets retain the declared finite-grid limitation.
 These are generating-point results, without a universal or real-discrepancy
 coverage guarantee.}
 \label{fig:coverage-delay}
\end{figure}

\begin{figure}[p]
 \centering
 \includegraphics[width=\linewidth]{coverage_delay_width.pdf}
 \caption{Median delay-interval widths depend on the uncertainty procedure and
 true serial dependence. Models, methods and simulation denominators match
 Fig.~\ref{fig:coverage-delay}. All panels share the same width scale. Profile
 widths sum accepted component lengths, excluding intervening gaps; the
 one-stage profile includes both pole-label modes. These width summaries do
 not establish nominal coverage. Exact counts, width distributions, exploration
 flags and interpolation sensitivity remain in the saved coverage summary.}
 \label{fig:coverage-width}
\end{figure}
